%% project-physics.tex
%%
%% ProjectionPhysics: Physics as Kinematics of a Flowing Space
%% Single-column preprint (REVTeX 4.2, APS, PRD)
%%
%% Compile:  tectonic projection-physics.tex   (or latex+bibtex+latex+latex)

\documentclass[aps,prd,preprint,superscriptaddress,amsmath,amssymb]{revtex4-2}

\usepackage{graphicx}   % Include figure files
\usepackage{bm}         % bold math
\usepackage{dcolumn}    % align table columns on decimal point
\usepackage{amsthm}
\newtheorem{postulate}{Postulate}[section]

\begin{document}

\title{Physics as Kinematics of a Flowing Space:\\ Mass Anchoring, Vector Light Speed, and Electrodynamics from a Single Space Field}

\author{Yuanjie Liu}
\email{yuanjieliu65@gmail.com}
\affiliation{Zhejiang University Binjiang Research Institute, Hangzhou, China}

\date{\today}

\begin{abstract}
We present a postulate-based framework in which space itself is in motion,
and all physics is the kinematics of a single space field $\bm{C}$ of fixed
magnitude $|\bm{C}|=c$.  The framework rests on four postulates: (P1) every
spatial point carries a velocity vector $\bm{C}$ with $|\bm{C}|^2=c^2$
(vector light speed --- $c$ is a property of space, not a speed limit for
matter); (P2) photons are exactly comoving with the flow, so they carry no
``anchoring'' and are massless; (P3) mass is the anchoring effect of internal
motion (spin) against the motion of space, so nonzero spin implies nonzero
mass; and (P4) space moves in three directions, the normal direction emerging
algebraically as $\sigma_3=i\sigma_1\sigma_2$, which realizes the
three-color/three-gluon structure of $\mathrm{SU}(3)$.  Under these
postulates the electromagnetic fields are kinematic derivatives of the space
field, $E=-\partial_t C$ and $B=\partial_x C$; Faraday's law becomes an
identity, Ampere--Maxwell becomes equivalent to the wave equation of $C$, and
the gauge redundancy of the vector potential is eliminated by the physical
condition $|\bm{C}|=c$.  Inhomogeneous flow yields the Gordon metric with
$\Phi=\frac12 v^2$ in the weak-field limit, reproducing the Newtonian limit
of general relativity; black holes appear as sinks of the flow whose horizon
is the light-speed surface, so that inescapability is a corollary of the
postulates rather than a dynamical result.  The three-direction hypothesis
predicts a glueball mass pattern $m=\sqrt{N}\,M_0$ that is consistent with
lattice QCD ($0^{++}$ at $\sqrt3\,M_0$ with $M_0\approx1$ GeV, $0.0\%$
deviation), and the mode counts $N=3,6,7$ align with the degeneracies
$d(n)=\frac12(n+1)(n+2)$ of the isotropic three-dimensional harmonic
oscillator.  All central theorems are formalized in the Lean 4 theorem prover
(no axioms, no admitted proofs).  Four further explorations are
formalized and checked against data: spin as the emergence of the
three-direction structure (SFS1--SFS5), double-slit interference as
interference of space-helix harmonics (DS1--DS4), the fractal cosmos
with the KBC void resolving the Hubble tension ($\delta^*=-0.249$
inside the observed range), and glueball three-direction coupling with
a massification gradient; the glueball $\sqrt{N}\,M_0$ sequence with
$N=3,6,7$ falls inside the lattice ranges for all three channels.  We state the epistemic status honestly:
at the present stage the framework is a self-consistent \emph{conceptual
reconstruction} --- mathematically and numerically indistinguishable from
standard physics wherever it overlaps --- and we catalogue the falsifiable
predictions and open problems that would elevate it beyond reinterpretation.
\end{abstract}

\pacs{04.20.-q, 03.50.De, 04.70.-s, 12.39.Mk, 03.65.Fd}

\maketitle

\section{Introduction}\label{sec:intro}

Hilbert's sixth problem asks for an axiomatic derivation of physics
\cite{Hilbert1900}.  The present project --- ProjectionPhysics, the
physical extension of the Hidden-space Bridge System (HIBS) of Liu and Xu
\cite{LiuXu} --- pursues a specific instance of this program.  HIBS starts
from a hidden space $S$ of ``hidden numbers'' and a non-injective projection
$\pi:S\to \mathbb{R}\oplus i\mathbb{R}$; all observables are functions of the
image $\pi(S)$ and the kernel $\ker\pi$, and no third degree of freedom
exists.  The classical complex plane $\mathbb{C}$ is not fundamental; it is
the observable slice of $S$ under $\pi$.

ProjectionPhysics adds a physical layer to this algebraic skeleton, guided
by a single intuition (Liu, 2026): \emph{space itself moves}, and light
speed is the equivalent speed of that motion.  From this intuition four
postulates are distilled (Sec.~\ref{sec:postulates}), and from them three
bodies of physics are derived:

\begin{enumerate}
  \item \textbf{Mass.}  Mass is not the excitation energy of a field; it is
  the \emph{anchoring} of internal motion (spin) against the motion of
  space.  Nonzero spin implies nonzero mass; exact comotion with the flow
  implies zero mass (Sec.~\ref{sec:mass}).
  \item \textbf{Electrodynamics.}  The electromagnetic fields are
  \emph{kinematic derivatives} of the space field $C$:
  $E=-\partial_t C$, $B=\nabla\times C$.  Maxwell's equations reduce to the
  wave equation of $C$; Faraday's law is an identity, and the gauge
  redundancy of the vector potential is killed by the physical condition
  $|\bm{C}|=c$ (Sec.~\ref{sec:em}).
  \item \textbf{Gravitation.}  Inhomogeneous flow distorts the metric; the
  Gordon metric follows, the weak-field potential is $\Phi=\frac12 v^2$,
  and black holes are sinks of the flow whose horizon is the light-speed
  surface (Sec.~\ref{sec:gravity}).
\end{enumerate}

Throughout, a methodological commitment is observed that is unusual for a
theory paper: \emph{every} central theorem is accompanied by a machine-checked
proof in the Lean 4 theorem prover \cite{deMoura2015,mathlib2020}, and every
numerical claim by a reproducible script in the public repository
\cite{ProjectionPhysics}.  This is not decoration: in a framework whose
content is mostly at the level of algebraic identities, formal verification
is what separates ``claimed'' from ``established''.

We also commit, equally unusually, to an explicit \emph{honesty contract}
(Sec.~\ref{sec:status}): a four-layer assessment of what each result
actually is --- mathematical identity, structural correspondence with known
physics, numerical matching, or conceptual reinterpretation --- and we
record, in writing, which real-discovery criteria are \emph{not} satisfied.
The reader will find the framework self-consistent, partially novel in its
derivation route, and currently indistinguishable from standard physics at
the level of formulas and numbers.

\section{Postulates of the flowing-space framework}\label{sec:postulates}

\subsection{Vector light speed (SLS1--SLS3)}

\begin{postulate}[Vector light speed]
\label{post:sls1}
There exists a vector field $\bm{C}(x)$ on space with
\begin{equation}
|\bm{C}(x)|^2 = c^2 \qquad \forall x,
\end{equation}
where $c$ is the universal speed of light.  $c$ is therefore a property of
\emph{space itself} --- the equivalent speed of space's own motion --- not a
maximum speed for matter \emph{in} space \cite{ProjectionPhysics-SLS}.
\end{postulate}

This is a genuine shift of ontology.  In the standard view, $c$ bounds the
propagation of signals through space; here, $c$ describes the motion of
space, and matter moves \emph{relative} to it.  The relative velocity of
matter with respect to the flow is
\begin{equation}
u = \frac{dx}{dt} - \bm{C}(x).
\end{equation}

\begin{postulate}[Photons are comoving]
\label{post:sls2}
Photons are exactly comoving with the flow: $u=0$.  A photon's worldline is
a streamline of $\bm{C}$.
\end{postulate}

\begin{postulate}[Mass is anchoring]
\label{post:mc1}
The mass of a particle is the \emph{anchoring effect} of its internal motion
against the motion of space: $m$ is a norm of the spinor flow $\sigma\psi$
generated by the internal state,
\begin{equation}
m \;:=\; \|\sigma\psi\| \geq 0,
\end{equation}
with $m=0$ iff the internal state is trivial.  Equivalently, in the hidden
number representation, $m^2 = h^2$ for $h\in\ker\pi$.
\end{postulate}

\begin{postulate}[Three-direction motion]
\label{post:three}
Space moves in three directions: planar motion (two directions) plus motion
along the normal vector, which is \emph{not} an independent input but
emerges algebraically from the planar generators,
\begin{equation}
\sigma_3 \;=\; i\,\sigma_1\sigma_2,
\qquad
(\sigma_1\sigma_2)^2 = -1,
\qquad
\sigma_1\sigma_2 + \sigma_2\sigma_1 = 0 .
\end{equation}
The three directions $(x,y,z)$ of space are identified with the three color
directions of $\mathrm{SU}(3)$: the normal vector $(1,1,1)$ is the
hypercharge direction, and the color-singlet profile $(c_0,c_1,c_2)=(1,1,1)$
lies along it.
\end{postulate}

Postulate~\ref{post:sls2} together with Postulate~\ref{post:sls1} implies
that a photon worldline satisfies $|dx/dt|=|\bm{C}|=c$, and hence the
proper time
\begin{equation}
d\tau^2 = dt^2 - \frac{dx^2}{c^2}
\end{equation}
vanishes identically along it: \emph{photons do not experience time because
they are the motion of space itself} (SM1).  Massive particles deviate from
the flow ($u\neq 0$), which is precisely their mass (SM3c): mass is
\emph{the inability to comove with space at the equivalent speed $c$}.
Time itself becomes a measure of deviation from the flow:
$d\tau^2>0$ iff $|dx|<c\,dt$ (SM2).

\subsection{Relation to the HIBS axioms}

The postulates are not independent of the HIBS algebra.  Postulate~3
interprets the HIBS kernel $\ker\pi$ as the carrier of internal (hidden)
motion; the Clifford generators $\sigma_i$ are the representation of the
projection algebra, and $i$ itself emerges from the anticommutation of the
planar generators, $(\sigma_1\sigma_2)^2=-1$, rather than being an axiom.
The Lean formalization makes this explicit: in \texttt{PauliMathlib.lean}
and \texttt{SpaceLightSpeed.lean} the statements
\texttt{normal\_direction\_emerges\_from\_plane} and
\texttt{planar\_motion\_products\_give\_i} are proved from the
anticommutation relations alone \cite{ProjectionPhysics-SLS}.

\section{Mathematical structure}\label{sec:math}

\subsection{The algebraic skeleton}

The HIBS projection structure provides the vocabulary: a non-injective
projection $\pi:S\to V$ with kernel $\ker\pi$, a quadratic form $Q$ on the
image, and a Clifford representation $\sigma$ of the form.  All theorems in
this paper live on top of this skeleton; the Lean files
\texttt{Definitions.lean}, \texttt{Kernel.lean}, \texttt{Completeness.lean},
and \texttt{Mass.lean} establish the elementary facts (projection pairs,
null cones, information conservation) \cite{ProjectionPhysics}.

\subsection{Spinor flow and anchoring}

Let $\psi\in\mathbb{C}^2$ be a two-component spinor and $\sigma_i$ the Pauli
matrices.  The \emph{spinor flow} in direction $i$ is $\sigma_i\psi$.  The
anchoring mass candidate is
\begin{equation}
m \;=\; \sum_i |(\sigma_i\psi)_0|^2 + |(\sigma_i\psi)_1|^2 ,
\end{equation}
and the core theorems (MC1, Lean) state
\begin{equation}
\psi \neq 0 \;\Longrightarrow\; m > 0 ,
\qquad
\psi = 0 \;\Longrightarrow\; m = 0 .
\end{equation}
The glueball minimal version (MC2) is a color-singlet triplet
$G=(a,b,c)$ with profile $(1,1,1)$:
\begin{equation}
m_G^2 = |a|^2 + |b|^2 + |c|^2 > 0
\qquad\text{whenever}\qquad (a,b,c)\neq(0,0,0).
\end{equation}
The number three is not accidental: the bridge theorem\footnote{Lean 4
theorem \texttt{three\_direction\_three\_glueball\_bridge}, module SLS5 of
\texttt{SpaceLightSpeed.lean}.} (SLS5) equates the norm squared of the
three-direction space motion (3) with the mass squared of the three-gluon
singlet (3).

\subsection{Spin-statistics and compositeness constraints}

An exploration module \texttt{CliffordSix.lean} and
\texttt{SpinStatistics.lean} verifies the well-known obstruction that pure
glueballs cannot assemble into fermions (spin-statistics), and that the
eight-dimensional Clifford spinor of $\mathrm{C}\ell(6)$ realizes the color
octet, with the branching $\mathrm{SU}(3)\supset\mathrm{SU}(2)$,
$3\to 2\oplus 1$ \cite{ProjectionPhysics}.  These modules are kept separate
from the main line because their physical content is standard.

\section{Mass: anchoring of internal motion}\label{sec:mass}

\subsection{The electron and the photon}

The framework assigns:
\begin{itemize}
  \item \textbf{Photon:} no internal structure, exactly comoving
  ($u=0$) $\Longrightarrow$ zero anchoring $\Longrightarrow$ $m=0$.
  \item \textbf{Electron:} spin means the normal-direction component
  $\sigma_3\psi\neq 0$ (the normal motion trajectory of the spin)
  $\Longrightarrow$ positive anchoring $\Longrightarrow$ $m>0$.
\end{itemize}
The mechanism chain is fully formalized (MC5$'$, SM3c, SM3d in Lean):
\begin{equation}
\psi\neq 0 \;\Longrightarrow\; \sigma_3\psi\neq 0
\;\Longrightarrow\; m>0
\;\Longrightarrow\; dx \neq c\,dt
\;\Longrightarrow\; d\tau^2 > 0 .
\end{equation}

A structural picture of the electron as a local flow structure (``tornado'')
is explored numerically: the electron is a swirl (curl $\nabla\times\bm{C}
\neq 0$, providing the anchoring mass) combined with a sink/source
(divergence $\nabla\cdot\bm{C}\neq 0$, providing the charge and the Coulomb
$1/r^2$ field), while the photon is a propagating wave of the same flow.
The classical self-energy of a point charge cut off at the flow scale
$r_0=0.20$\,fm evaluates to $U(r_0)=3.60$\,MeV $\approx 7.05\,m_e$, and at
the classical electron radius $r_e=2.82$\,fm to $U(r_e)=0.255$\,MeV
$\approx 0.5\,m_e$; the electron mass lies between the two estimates
\cite{ProjectionPhysics-electron}.  This is reported honestly as an order-of-
magnitude coincidence, not a derivation: $r_0$ is a fitted input, and the
classical $4/3$ and $1/2$ factor problems remain.

\subsection{Glueball spectrum: $\sqrt{N}\,M_0$}

The three-direction hypothesis makes a quantitative statement about the
lightest glueball.  With $M_0$ the mass scale of one spatial mode,
\begin{equation}
m_G(N) = \sqrt{N}\,M_0 ,
\end{equation}
the lattice values of \cite{MorningstarPeardon1999,Chen2006} are matched as
follows:

\begin{table}[h]
\caption{Glueball masses from the $\sqrt{N}\,M_0$ rule with $M_0=0.987$\,GeV
($r_0=\hbar c/M_0=0.200$\,fm), compared with lattice QCD
\cite{MorningstarPeardon1999,Chen2006}.  The fit parameter is $M_0$; $N$ is
the mode count.}
\label{tab:glueball}
\begin{ruledtabular}
\begin{tabular}{lcccc}
Channel & $N$ & $\sqrt{N}\,M_0$ (GeV) & Lattice (GeV) & Deviation \\
\hline
$0^{++}$ & 3 & 1.710 & $1.71\pm0.05$ & $0.0\%$ \\
$2^{++}$ & 6 & 2.418 & $2.40\pm0.12$ & $0.8\%$ \\
$0^{-+}$ & 7 & 2.612 & $2.56\pm0.15$ & $2.0\%$ \\
\end{tabular}
\end{ruledtabular}
\end{table}

The primary match is the $0^{++}$ glueball: $\sqrt{3}\,M_0=1.710$\,GeV
against the lattice value $1.71\pm0.05$\,GeV, a $0.0\%$ deviation with
$M_0\approx 1$\,GeV --- the natural nuclear scale --- and
$r_0\approx0.2$\,fm, the proton Compton wavelength scale.

\subsection{Why $N=3,6,7$: harmonic-oscillator degeneracies}

The mode counts admit a first-principles candidate.  The degeneracies of the
isotropic three-dimensional harmonic oscillator are
\begin{equation}
d(n) = \frac{(n+1)(n+2)}{2} = 1,\,3,\,6,\,10,\,15,\,21,\,\ldots
\end{equation}
and the lattice mode counts align:
\begin{equation}
0^{++}\to N=3=d(1), \qquad
2^{++}\to N=6=d(2), \qquad
0^{-+}\to N=7=d(3)-d(1).
\end{equation}
Against the two-dimensional alternative $N=2,4,5$, the three-direction
interpretation wins by $2.75\sigma$ in the weighted deviation
($\sigma_A=0.263$ vs.\ $\sigma_B=3.012$).  The physical picture is a glueball
as a bound quantum of the three-direction space flow --- a standing-wave
combination of metric perturbations, with $m^2\propto$ (number of modes)
$\times M_0^2$.  Because the flow is geometric (volume-preserving,
$\det g=-1/c^2$) rather than material, no Michelson--Morley-type anisotropy
is induced.

Two honest caveats are recorded: (i) the rule $0^{-+}\to d(3)-d(1)$ is an
ad-hoc symmetry subtraction (removing the modes with $0^{++}$ symmetry) and
lacks a symmetry derivation; (ii) $M_0$ itself has no first-principles
origin (it is the one fit parameter).  The predictions that follow from the
rule are listed in Sec.~\ref{sec:predictions}.

\section{Electrodynamics as kinematics of the space field}\label{sec:em}

\subsection{The central derivation}

Let $C$ denote the space field of Postulate~\ref{post:sls1}.  Define the
electromagnetic fields as \emph{kinematic derivatives} of $C$
\cite{ProjectionPhysics-MS}:
\begin{equation}
\label{eq:kinematic}
E \;=\; -\partial_t C , \qquad
B \;=\; \partial_x C \quad (1{+}1\mathrm{D}), \qquad
B \;=\; \nabla\times C \quad (3\mathrm{D}).
\end{equation}
Then the following are immediate and formalized (MS1--MS5, Lean, zero
admitted axioms):

\begin{enumerate}
  \item \textbf{Faraday's law is an identity} (MS2).  With
  $E=-\partial_t C$ and $B=\partial_x C$,
  \begin{equation}
  \partial_t B = \partial_t\partial_x C = \partial_x\partial_t C = -\partial_x E ,
  \end{equation}
  where the middle step is the commutation of mixed partials (MS1).  Faraday
  is not an independent law; it is the statement that $C$ is a single field.
  In 3D the same argument gives $\partial_t B = -\nabla\times E$, and
  $\nabla\cdot B = \nabla\cdot(\nabla\times C)=0$ is automatic.

  \item \textbf{Ampere--Maxwell is the wave equation of $C$} (MS3).
  \begin{equation}
  \partial_t E = -c^2\,\partial_x B
  \;\;\Longleftrightarrow\;\;
  \partial_t^2 C = c^2\,\partial_x^2 C .
  \end{equation}
  The two sides are literally the same equation.  Hence the independent
  content of the vacuum Maxwell system is a single wave equation for $C$;
  $E$ and $B$ are its apparent forms.

  \item \textbf{Sources are the driving term} (MS5).  A charge/current
  $J$ enters as a source:
  \begin{equation}
  \partial_t^2 C = c^2\,\partial_x^2 C + J ,
  \end{equation}
  so charge is not an external ``bullet'' but the driving term of the space
  field equation itself.  In the static limit, $C''\approx -J$, verified
  numerically.

  \item \textbf{Gauge redundancy is killed by physics} (MS4).  In the
  standard vector-potential formulation, $C$ would be the gauge-dependent
  potential $A$.  Here $C$ is physically fixed by
  Postulate~\ref{post:sls1} ($|\bm{C}|=c$ everywhere): the gauge freedom is
  eliminated not by a choice of convention but by the physical condition
  that $C$ is the actual velocity field of space.
\end{enumerate}

In short: \emph{electromagnetism is the kinematics of the space field}.
Maxwell's equations do not describe two fields with two independent laws;
they describe one field with one wave equation, and the familiar
electromagnetic laws are consequences of the definition of $E$ and $B$ as
derivatives of $C$.

\subsection{Consistency with the flow postulates}

The identification $c=1/\sqrt{\mu_0\varepsilon_0}$ (MF2) is preserved: $c$
is simultaneously the flow speed of space and the propagation speed of the
$C$ wave.  Dispersion $\omega=\pm ck$ holds as a factorization
$(\omega-ck)(\omega+ck)=0$ (MF1), and the photon-comoving postulate (SLS2)
implies $|v_{\mathrm{photon}}|=|C|=c$, so $d\tau^2=0$ identically (MF4,
self-consistent).  A leapfrog FDTD simulation with CFL$=1$ propagates the
$C$ wave packet at exactly $1.000\,c$, and the Faraday and Ampere residuals
are zero at machine precision for both analytic traveling waves and the
numerical leapfrog solution \cite{ProjectionPhysics-MS}.

\subsection{What is genuinely new here}

The kinematic content --- that Maxwell's equations are equivalent to the
wave equation of the vector potential --- is standard electromagnetic theory.
What is repository-specific is the \emph{physicalization} of the potential:
$C$ is not an auxiliary field but the velocity field of space, with the
fixed magnitude $|\bm{C}|=c$ that eliminates gauge freedom.  This is a
conceptual step; its falsifiability is discussed in Sec.~\ref{sec:status}.

\section{Gravitation and strong-field structure}\label{sec:gravity}

\subsection{From momentum conservation to the Gordon metric}

The route from the postulates to gravity follows the classical Gordon
construction \cite{Gordon1923}: an inhomogeneous flow $v(x)$ distorts the
metric.  With $g=\mathrm{diag}(1,-1/c^2)$ in the comoving frame, a flow
along $x$ gives the Gordon metric (SG1)
\begin{equation}
g_{\mu\nu} = \begin{pmatrix} 1-v^2/c^2 & v/c^2 \\ v/c^2 & -1/c^2 \end{pmatrix},
\qquad \det g = -\frac{1}{c^2} ,
\end{equation}
and the proper time (SG8)
\begin{equation}
d\tau^2 = dt^2 - \frac{(dx-v\,dt)^2}{c^2} .
\end{equation}
The key quantitative result (SG11) is the weak-field match: with
$g_{tt}=1-v^2/c^2$ from the flow and $g_{tt}=1-2\Phi/c^2$ from standard weak
field general relativity,
\begin{equation}
\Phi \;=\; \frac12 v^2 ,
\end{equation}
so the Newtonian acceleration is
\begin{equation}
a = -\nabla\left(\tfrac12 v^2\right) = -v\,\nabla v ,
\end{equation}
i.e., gravity is the inhomogeneity of the flow (SLS6).  The full chain
--- inhomogeneous flow $\to$ Gordon metric $\to$ Christoffel symbols $\to$
geodesic equation $\to$ Newtonian limit --- is verified numerically with the
generic formula $\Gamma^\lambda_{\mu\nu}=\frac12 g^{\lambda\sigma}(\partial_\mu
g_{\sigma\nu}+\partial_\nu g_{\sigma\mu}-\partial_\sigma g_{\mu\nu})$,
including a correction pass that recomputed all eight Christoffel components
\cite{ProjectionPhysics-GR}.  Proper time remains positive for massive
particles ($d\tau^2>0$) and vanishes for photons, whose worldlines satisfy
$|dx-v\,dt|=c\,dt$ \emph{relative to the flow} --- the correct photon
condition, as opposed to the comoving condition $dx=v\,dt$, which would
incorrectly give $d\tau^2=dt^2$ (a modeling pitfall recorded in the
formalization) \cite{ProjectionPhysics-GR}.

\subsection{Black holes as sinks of the flow}

The postulates (rather than a field equation) determine the strong-field
picture.  A black hole is a region where the flow converges to a point:
$C=-c\,\hat{r}$ inside $r<r_h$ (a sink), with all streamlines terminating at
the singularity \cite{ProjectionPhysics-MF}.  The horizon is the
light-speed surface (BH1):
\begin{equation}
g_{tt}=0 \;\Longleftrightarrow\; |v|=c .
\end{equation}
Inescapability is then a \emph{corollary of Postulate~\ref{post:sls2}}:
photons are streamlines, and escaping would require moving against the flow,
which contradicts exact comotion (PH2).  No appeal to general relativity is
needed.  Numerically, $|C|=c$ holds to $<10^{-12}$ along all streamlines,
$d\tau^2=0$ to $<10^{-12}$, and an outward photon freezes at the horizon
($dx/dt=0$), the coordinate picture matching Schwarzschild.  The white hole
is the reversed flow (BH5--BH6: $g_{tt}$ invariant, $g_{tx}$ flips), and a
wormhole is a sub-light flow tube whose throat has a peak speed
$v_{\mathrm{peak}}<c$ (WH1: at the throat $v=0$ and the light cone is
restored), through which a photon passes with $d\tau=0$ throughout
\cite{ProjectionPhysics-BH}.  This is structurally identical to the
Hamilton--Lisle river model \cite{HamiltonLisle2008}.

\section{Numerical and computational validation}\label{sec:validation}

All numbers in this paper are produced by scripts in the repository
\cite{ProjectionPhysics}.  Table~\ref{tab:validation} collects the
verification results.

\begin{table}[h]
\caption{Computational validation summary.  ``Residual zero'' means
machine-precision zero.}
\label{tab:validation}
\begin{ruledtabular}
\begin{tabular}{ll}
Check & Result \\
\hline
$C$-wave packet speed (leapfrog FDTD, CFL$=1$) & $1.000\,c$ \\
Faraday residual (analytic / leapfrog) & $0.0$ / $0.0$ \\
Ampere--Maxwell residual (analytic / leapfrog) & $0.0$ / $0.0$ \\
Static charge drives $C$ (MS5) & $C''\approx -J$, residual $0.033$ \\
Glueball $0^{++}$ ($\sqrt3\,M_0$) & $1.710$ vs.\ $1.71\pm0.05$ GeV ($0.0\%$) \\
Glueball $2^{++}$ ($\sqrt6\,M_0$) & $2.418$ vs.\ $2.40\pm0.12$ GeV ($0.8\%$) \\
Glueball $0^{-+}$ ($\sqrt7\,M_0$) & $2.612$ vs.\ $2.56\pm0.15$ GeV ($2.0\%$) \\
Horizon condition $g_{tt}=0\Leftrightarrow|v|=c$ & exact \\
Outward photon at horizon & $dx/dt=0$ (frozen) \\
Time dilation at horizon & $d\tau/dt\to 0$ \\
Wormhole throat peak & $v=0.8\,c<c$, no horizon \\
$d\tau^2$ along streamlines & $<10^{-12}$ \\
$|\bm{C}|=c$ along streamlines & $<10^{-12}$ \\
CHSH: helix model $S$ & $2.0005\pm0.004$ (local bound $2$) \\
CHSH: quantum prediction & $2\sqrt2=2.8284$; Aspect 1982: $2.697$ \\
Electron self-energy $U(r_0=0.20\,\mathrm{fm})$ & $3.60$ MeV $=7.05\,m_e$ \\
Electron self-energy $U(r_e=2.82\,\mathrm{fm})$ & $0.255$ MeV $=0.50\,m_e$ \\
\end{tabular}
\end{ruledtabular}
\end{table}

\subsection{Lean formalization}

The theorem prover component is summarized in Table~\ref{tab:lean}.  Every
listed module compiles with zero \texttt{sorry} (no admitted axioms) and
zero warnings, under Lean 4.

\begin{table}[h]
\caption{Lean 4 formalization summary.  Job counts are from
\texttt{lake build} on the current revision.}
\label{tab:lean}
\begin{ruledtabular}
\begin{tabular}{ll}
Module (file) & Theorems \\
\hline
\texttt{SpaceLightSpeed.lean} & SLS1--SLS5: universal $|\bm{C}|=c$; photon zero anchoring; \\
 & electron positive anchoring; $\sigma_3=i\sigma_1\sigma_2$ emergence; \\
 & three-direction/three-gluon bridge \\
\texttt{MinimalCore.lean} (Archive) & MC1--MC2, MC1h: nonzero spin $\Rightarrow$ nonzero mass; \\
 & $m_G^2=|a|^2+|b|^2+|c|^2$ \\
\texttt{SpaceMetric.lean} & SM1--SM6: photon $d\tau=0$; mass $\Leftrightarrow$ deviation \\
 & from flow; metric consistency; $\det g=-1/c^2$ \\
\texttt{SpaceGravity.lean} & SG1--SG11: Gordon metric; $\Phi=\frac12 v^2$; Newtonian limit \\
\texttt{MaxwellSpace.lean} & MS1--MS5: Faraday identity; Ampere $\Leftrightarrow$ wave equation; \\
\texttt{(Explorations)} & sources as driving terms (4131 jobs, zero sorry) \\
\texttt{MaxwellFlow.lean} & MF1--MF6, PH1--PH2: dispersion; $c=1/\sqrt{\mu_0\varepsilon_0}$; \\
\texttt{(Explorations)} & inescapability as postulate corollary (4130 jobs) \\
\texttt{BlackHoleWormhole.lean} & BH1--BH7, WH1: horizon $=$ light-speed surface; white \\
\texttt{(Explorations)} & hole; wormhole throat (4128 jobs) \\
\texttt{LorentzRebuild.lean} & LR1--LR5: boost preserves $\eta=\mathrm{diag}(-1,1)$; \\
 & $\gamma^2-\gamma^2\beta^2=1$; velocity addition; $\beta<1$ \\
\texttt{DiracMathlib.lean} & DB1$'$--DB6$'$: $(\gamma^0)^2=1$, $(\gamma^i)^2=-1$; mass $=$ \\
 & chiral coupling \\
\texttt{EntanglementHelix.lean} & EH1--EH4: CHSH local bound $|S|\leq2$ (Bell kernel) \\
\texttt{(Explorations)} & \\
\end{tabular}
\end{ruledtabular}
\end{table}


\section{Recent explorations: spin, double-slit, fractal cosmos, and glueball dynamics}\label{sec:explorations}

Four further checks of the framework were carried out after the first
version of this paper; all are formalized in Lean (zero \texttt{sorry})
and validated numerically (artifacts and scripts in the repository).

\subsection{Spin as the emergence of the three-direction structure}

Within the three-direction postulate (P4), the spin structure is not an
input: the three directions generate a Clifford algebra, and the imaginary
unit emerges as the oriented volume element
\begin{equation}
i = \sigma_1\sigma_2\sigma_3 , \qquad
[\sigma_i,\sigma_j] = 2i\varepsilon_{ijk}\sigma_k , \qquad
S^2 = \tfrac34\,\hbar^2 = s(s{+}1)\,\hbar^2 \quad (s=\tfrac12) .
\end{equation}
The value $s=\tfrac12$ is the Casimir value forced by the $2$-dimensional
(irreducible) representation of $C\ell(3)$, and the sign flip under a
$2\pi$ rotation (``fermionic'' behaviour) is the double-cover geometry of
$\mathrm{SU}(2)$ over $\mathrm{SO}(3)$ --- not a quantum postulate.  This
answers, within the framework, the question Dirac left open (where does the
``complex intrinsic property'' come from? --- it is the orientation of the
three directions) while honestly recording that the choice of the
2-dimensional representation and the value of $\hbar$ remain inputs
(``second input'' gap).

\subsection{Double-slit interference as interference of space-helix harmonics}

Photons are comoving with the flow (P2); if the flow carries a helical
structure (circular motion in the plane plus translation along the normal,
the two planar modes of P4), then the projection of the helix onto any
plane is a harmonic wave ($x=R\cos\omega t$; the $xy$ projection is exactly
a circle, $x^2+y^2=R^2$, DS1).  Interference fringes are then the algebraic
consequence of adding two harmonics,
\begin{equation}
\cos\alpha+\cos\beta=2\cos\!\tfrac{\alpha+\beta}{2}\,\cos\!\tfrac{\alpha-\beta}{2},
\qquad I(\theta)=4I_0\cos^2\!\big(\tfrac{\pi d\sin\theta}{\lambda}\big),
\end{equation}
and observation (which-slit information) removes the interference cross
term $2\cos\Delta\varphi$ in $|z_1+z_2|^2$ --- the ``compression of the
space harmonics'' --- leaving two bands (DS4).  The mathematics is
identical to standard wave optics and decoherence; the re-interpretation
is conceptual (the wave is the projection of the space helix, the particle
is the compressed-harmonic localisation).

\subsection{Fractal cosmos, the KBC void, and the Hubble tension}

If cosmic structure is fractal (observed galaxy distribution
$D\simeq1.2{-}1.5$), then in the flow picture the expansion is the
divergence of the flow, $H(\bm{x})=H_0\big(1-\delta(\bm{x})/3\big)$
(top-hat linear theory), and voids are the necessary empty regions of the
fractal ``skeleton'' (cytoskeleton analogy).  The local Hubble constant
measured inside a deep void is elevated: reproducing the observed tension
$\Delta H/H=(73.0-67.4)/67.4=8.3\%$ requires a void depth
\begin{equation}
\delta^* = -3\,\Delta H/H = -0.249 ,
\end{equation}
which falls inside the observed KBC void range
$\delta\in[-0.3,-0.2]$; a simulated fractal flow field (96$^3$ FFT,
Poisson-solved, spectral divergence identity exact to machine precision)
produces a largest void with central $\delta=-0.27$ and an $9\%$ Hubble
boost, consistent with the $5.6$ km/s/Mpc tension.  This reproduces the
known KBC-void resolution of the tension (Keenan et al.\ 2013, Shanks et
al.\ 2019) in the flow language; it is a consistency check, not a new
prediction.

\subsection{Glueball three-direction coupling and the massification gradient}

Describing the glueball in the three-direction version of the modified
Maxwell system ($V=g\,C_1C_2C_3$ three-linear coupling) yields a stable
localised bound state whose mass grows monotonically with the coupling;
the mass-squared decomposes into the three-direction gradient sum
$m^2=\sum_i(\partial_x C_i)^2$ (equal gradients $\Rightarrow
m\propto\sqrt3\,g$, the same $\sqrt3$ factor as the mode rule), and the
``gradient that resists the space flow'' is the massification mechanism
($\Phi=\tfrac12 v^2$, SG11).  With the mode rule of
Sec.~\ref{sec:mass} ($N=3,6,7$), the updated lattice comparison
(Morningstar--Peardon 1999; Chen et al.\ 2006) gives

\begin{table}[h]
\caption{Glueball masses from the $\sqrt{N}\,M_0$ rule with $M_0=0.93$\,GeV
(central), updated with the three-direction coupling check.}
\label{tab:glueball2}
\begin{ruledtabular}
\begin{tabular}{lcccc}
Channel & $N$ & $\sqrt{N}\,M_0$ (GeV) & Lattice (GeV) & Verdict \\
\hline
$0^{++}$ & 3 & 1.61 & $1.475{-}1.75$ & inside \\
$2^{++}$ & 6 & 2.28 & $2.15{-}2.40$ & inside \\
$0^{-+}$ & 7 & 2.46 & $2.30{-}2.60$ & inside \\
\end{tabular}
\end{ruledtabular}
\end{table}

All three channels fall inside the lattice ranges (previously $2^{++}$ was
$3{-}13\%$ low with $N=5$; the oscillator degeneracy $N=6=d(2)$ brings it
inside).  The inverse-derived $M_0\simeq0.93\pm0.07$\,GeV is consistent
across the three channels to better than $7\%$.  The $N$-sequence and the
$M_0$ calibration remain fits (second-input gap), and the coupling model is
a $1{+}1$-dimensional toy, so the verdict is: data consistent at the
order-of-magnitude level; first-principles prediction not yet achieved.

\subsection{Information as symplectic volume, photon helices, and counting the strands (2026-08-16)}

A fast-moving exploration line (formalized in
\texttt{ProjectionPhysics/Explorations/QFTFlow.lean}, 4146 jobs, zero
`sorry`; numerical checks N1--N28 at machine precision) extends the
framework in four directions.

\textbf{Information as symplectic volume.} Unitary transport (flow
propagation) preserves $\det(AA^\dagger)$ (GQC1): information, defined as
the symplectic correlation volume of the twistors, is conserved by the
flow. Causality is lattice locality (GQC2): a nearest-neighbour jump
matrix has bandwidth $t$ after $t$ steps, so the propagation kernel
vanishes outside the lattice light cone. In a non-uniform flow the
effective speed is $1+v_{\rm flow}$ (GQC3): superluminal \emph{equivalent}
motion is a geometric (coordinate-dragging) effect --- the
Alcubierre/Painlev\'e--Gullstrand structure --- while signals never exceed
the local speed of light. No new physics is claimed; these belong to the
explanation layer.

\textbf{Flow momentum and the four forces.} With the momentum postulate
$p = m(C-v)$ ($C$ the vector light speed, $v$ the body velocity), the
Leibniz rule gives (GQF)
\[
\frac{d}{dt}\big[m(C-v)\big] = \dot m\,C + m\dot C - \dot m\,v - m\dot v,
\]
four channels identified with the electric ($\dot m C$), nuclear
($m\dot C$), magnetic ($\dot m v$) and gravitational/inertial ($m\dot v$)
forces. This is an algebraic identity; the coupling constants are not
derived.

\textbf{Photon = excited electron helix.} The photon is an excited
electron whose mass and charge have vanished; the helix (two symmetric
helices, model B) has vanishing azimuthal momentum, hence purely axial
momentum and $E^2=|p|^2$ (massless) --- the two models share the light
cone.

\textbf{Counting the strands (matching attempt).} The quantum relation
\[
E = \hbar\omega = \frac{h}{2\pi}\,2\pi f = h\,f
\]
is reinterpreted as \emph{angular momentum per turn times turns per
second} (GQR): $\hbar$ becomes a per-radian angular momentum, a countable
geometric quantity, and the turns are the (countable) frequency. The
numerical check is exact (relative error $0$). The glueball rule
$m=\sqrt{N}\,M_0$ reads the same way: $N=3$ strands (three anchored
directions); colour $=3$ is the same ``3'' (three spatial directions);
$8$ gluons $= 2^3 = \dim C\ell(6)$ (TW5). Honest limits: the quark mass
spectrum ($2.2$\,MeV $\to$ $173$\,GeV) shows no strand-counting pattern
(QCD free parameters); $L=h$ still rests on the measured photon spin; a
first-principles derivation of $\hbar$ from the helix geometry is the open
next step.

\section{Relation to existing physics}\label{sec:relation}

\subsection{What is standard, restated}

Where the framework touches existing physics, the formulas and numbers agree
with it: the Gordon metric is a classical construction
\cite{Gordon1923}; the river model of black holes is Hamilton and Lisle
\cite{HamiltonLisle2008}; Maxwell's equations as the wave equation of a
potential is textbook; $d\tau^2=0$ for photons is Minkowski; the glueball
spectrum is lattice QCD \cite{MorningstarPeardon1999,Chen2006}.  The
framework currently contains \emph{no} experimental number that standard
physics does not also produce.

\subsection{Four-layer assessment}\label{sec:status}

For each result we record its status on four layers
\cite{ProjectionPhysics-status}:

\begin{enumerate}
  \item \textbf{Mathematical identity.}  True by construction: e.g., the
  Faraday identity, Ampere $\Leftrightarrow$ wave equation, $g_{tt}=0
  \Leftrightarrow |v|=c$.  These are rigorous (machine-checked) but
  content-free as physics.
  \item \textbf{Structural correspondence.}  Known physics restated: the
  vector-potential formulation, the Gordon/river geometry, the
  spin-statistics obstruction.
  \item \textbf{Numerical matching.}  Consistent but not new: the glueball
  $\sqrt{N}M_0$ rule has one fit parameter $M_0$; the electron self-energy
  coincidence sits between $0.5\,m_e$ and $7\,m_e$ with a fitted cutoff.
  \item \textbf{Conceptual reinterpretation.}  Unfalsifiable at present:
  ``gravity is inhomogeneous flow'', ``mass is anchoring'', ``$c$ is the
  speed of space'' --- the explanation layer changes while the formulas do
  not.
\end{enumerate}

The one result that rises above layer 1 is the Maxwell-space derivation
(Sec.~\ref{sec:em}): although the algebraic kernel is the standard vector
potential formulation, the \emph{physicalization} of the potential
($|\bm{C}|=c$) and the reduction of Maxwell's system to a single wave
equation of a physically fixed field is repository-specific content.  Even
here, the honest statement is that it is a derivation-route novelty, not a
new prediction.

\section{Discussion}\label{sec:discussion}

\subsection{Predictions}\label{sec:predictions}

The framework makes the following falsifiable statements.

\begin{enumerate}
  \item \textbf{P1 --- inside the horizon, light inevitably falls into the
  singularity.}  This follows from the comoving postulate and is verified
  numerically \cite{ProjectionPhysics-MF}.
  \item \textbf{P2 --- gravitational redshift}
  \begin{equation}
  \frac{\omega_2}{\omega_1} = \frac{c-v_2}{c-v_1}
  \end{equation}
  with the flow profile $v(r)$ as the single free input.  The direction
  agrees with GR; numerical agreement with GR requires the flow profile to
  take the GR shape, which the repository does not derive from first
  principles.
  \item \textbf{P3 --- transverse electromagnetic modes vanish inside the
  horizon} (the field becomes purely radial) \cite{ProjectionPhysics-MF}.
  \item \textbf{P4 --- isotropy is preserved.}  Since $|\bm{C}|=c$ in every
  direction, light speed is isotropic by construction; this is structurally
  untestable, and we say so.
  \item \textbf{Glueball spectrum.}  The harmonic-oscillator mode rule
  predicts, for $n=4,5,6$,
  \begin{equation}
  m = \sqrt{15},\sqrt{21},\sqrt{28} \times M_0 = 3.82,\,4.52,\,5.22\,\mathrm{GeV},
  \end{equation}
  and, if the $d(n)-d(1)$ subtraction rule holds, $n=4$ gives
  $\sqrt{12}\,M_0=3.42$\,GeV.  Lattice data on $1^{++}$, $2^{-+}$, $3^{++}$
  channels can discriminate.
  \item \textbf{Entanglement under inhomogeneous flow.}  In the
  entanglement-helix exploration \cite{ProjectionPhysics-helix}, a
  gravitational (inhomogeneous-flow) phase shift would modify the
  correlation function $E(\Delta)$; the local bound $|S|\leq2$ is proven,
  and the helix geometry saturates it ($S=2.0005\pm0.004$), while the
  quantum value $2\sqrt2$ requires Born-rule readout --- an honest place
  where the framework must either produce Born's rule or fail.
\end{enumerate}

\subsection{Open problems}

The following gaps are structural, not cosmetic:

\begin{itemize}
  \item \textbf{The $966\times$ gap.}  The electron anchoring mechanism has
  the right direction (nonzero spin $\Rightarrow$ nonzero mass) but the
  magnitude, at the flow scale $r_0$, overshoots $m_e$ by a factor of
  order $10^3$; every candidate formula collapses to an identity
  ($m\equiv M_0/2$).  An independent input not of order $\hbar/r_0$ is
  needed.
  \item \textbf{Origin of $M_0$.}  The glueball fit parameter has no
  first-principles derivation; connecting it to a confinement scale requires
  QCD dynamics the repository does not have.
  \item \textbf{The $N$-sequence.}  $(3,6,7)$ is not yet derived; the
  $d(3)-d(1)$ subtraction for $0^{-+}$ is ad hoc.
  \item \textbf{Continuum 3D calculus.}  The 3D discrete curl/divergence
  identities are now formalized (SF1--SF5: $\operatorname{div}\!\operatorname{curl}=0$
  automatic, 3D Faraday automatic, gradient curl-free); the continuum
  version and continuity remain open.
  \item \textbf{The $N$-sequence.}  $(3,6,7)$ now has a dynamical candidate
  (three-linear coupling, GC1--GC4) but is still not derived; the
  $d(3)-d(1)$ subtraction for $0^{-+}$ remains ad hoc.
  \item \textbf{Representation choice and $\hbar$.}  Why the electron uses the
  2-dimensional representation, and the numerical value of $\hbar$, have no
  origin in the framework (second-input gap).
  \item \textbf{Missing standard content.}  Charge quantization, the value
  of $e$, QCD confinement, the complete curvature tensor and the Einstein
  field equations, and Born's rule are all outside the framework.
  \item \textbf{Lorentz violation?}  The only physical exit from pure
  reinterpretation is inhomogeneous flow inducing Lorentz violation; no
  experimental scheme or magnitude estimate exists yet.
\end{itemize}

\section{Conclusion}\label{sec:conclusion}

We have presented a postulate-based framework --- space in motion,
$|\bm{C}|=c$ --- from which mass, electrodynamics, and gravitation emerge as
kinematics.  The framework is internally consistent, every central claim is
machine-checked in Lean 4, and every number is reproducible from the
repository.  Its distinctive content is the physicalization of the vector
potential ($|\bm{C}|=c$) and the consequent reduction of Maxwell's equations
to the wave equation of a single space field, plus the three-direction
realization of $\mathrm{SU}(3)$ with its $\sqrt{N}M_0$ glueball pattern.

Its epistemic status is stated without inflation: at present this is a
\emph{conceptual reconstruction} of known physics --- rigorous, coherent,
and indistinguishable from standard physics in formulas and numbers.  Real
discovery criteria (first-principles derivation of $M_0$, resolution of the
$966\times$ gap, a new falsifiable prediction with an experimental scheme)
are not yet met.  The predictions of Sec.~\ref{sec:predictions} are the
tests that would change this verdict.

\begin{acknowledgments}
The author thanks X.~Xu for the HIBS collaboration, and the Lean community
for the formalization infrastructure on which the theorem-prover component
rests.
\end{acknowledgments}

\bibliography{projection-physics}

\appendix

\section{Lean formalization overview}\label{app:lean}

The formalization lives in the public repository
\texttt{github.com/logos-42/Hibs-Physics} (directory
\texttt{ProjectionPhysics/}).  Conventions: theorems are named by module
prefix (MS, MF, BH, SG, SM, SLS, MC, LR, EH); a prime denotes a mathlib
rewrite of an earlier core proof; \texttt{(Explorations)} modules are
newer-direction results kept separate from the main line.  Every file
compiles under \texttt{lake build} with zero \texttt{sorry} and zero
warnings; the job counts in Table~\ref{tab:lean} are from the current
revision.

\section{Proof sketches}\label{app:proofs}

\subsection{Faraday's law is an identity (MS2)}\label{app:faraday}

In $1{+}1$ dimensions, with $E=-\partial_t C$ and $B=\partial_x C$:
\[
\partial_t B \;=\; \partial_t\partial_x C
\;\overset{\mathrm{(MS1)}}{=}\;
\partial_x\partial_t C \;=\; -\partial_x E .
\]
In 3D, $B=\nabla\times C$ gives $\partial_t B=\nabla\times\partial_t C
=-\nabla\times E$, and $\nabla\cdot B=\nabla\cdot(\nabla\times C)=0$
identically.  Both directions are formalized; the commutation step (MS1) is
the only ingredient.

\subsection{Ampere--Maxwell is the wave equation (MS3)}\label{app:ampere}

Using the definitions,
\[
\partial_t E + c^2\,\partial_x B
\;=\; -\partial_t^2 C + c^2\,\partial_x^2 C ,
\]
so $\partial_t E = -c^2\,\partial_x B$ holds iff
$\partial_t^2 C = c^2\,\partial_x^2 C$.  The equivalence is a direct
substitution (Lean: MS3, both directions).

\subsection{Horizon $=$ light-speed surface (BH1)}\label{app:bh}

For the Gordon metric $g_{tt}=1-v^2/c^2$,
\[
g_{tt}=0 \;\Longleftrightarrow\; v^2=c^2 \;\Longleftrightarrow\; |v|=c .
\]
Inside the horizon $|v|>c$ gives $g_{tt}<0$ (BH7, role swap seed).

\subsection{Weak-field potential (SG11)}\label{app:sg11}

Equating the flow expression and the GR expression for $g_{tt}$,
\[
1-\frac{v^2}{c^2} = 1-\frac{2\Phi}{c^2}
\;\Longrightarrow\; \Phi = \frac{v^2}{2},
\]
and the geodesic equation yields $a=-\nabla\Phi=-v\,\nabla v$ (Newtonian
limit).

\subsection{CHSH local bound (EH1--EH2)}\label{app:chsh}

For dichotomic observables $A_a,B_b\in\{\pm1\}$,
\[
S = E(a,b)+E(a,b')+E(a',b)-E(a',b'),
\]
and $|S|\leq2$ follows from $A_aB_b+A_aB_{b'}+A_{a'}B_b-A_{a'}B_{b'}
=A_a(B_b+B_{b'})+A_{a'}(B_b-B_{b'})\in\{\pm2\}$, since each parenthesis is
$0$ or $\pm2$.  This is the algebraic kernel of Bell's theorem
\cite{Bell1964}; the helix geometry saturates the bound numerically
($S=2.0005\pm0.004$).

\subsection{Glueball mass positivity (MC2)}\label{app:mc2}

For $G=(a,b,c)$ with $m_G^2=|a|^2+|b|^2+|c|^2$, if $(a,b,c)\neq(0,0,0)$
then at least one $|a_i|^2>0$, hence $m_G^2>0$; conversely $m_G^2=0$ forces
$a=b=c=0$.  Trivial in content, but it is the formal anchor of ``three
directions give mass''.

\end{document}
